\chapter{Kafka trong cụm: Strimzi và luồng gửi mail}
\label{ch:strimzi}

Cột mốc~6 chuyển broker của Chương~10 từ Compose vào k3s và triển khai
service cuối cùng --- notification-service --- để luồng \emph{sự kiện}
đầu tiên chạy trọn vẹn trong cụm: đăng ký một người dùng, nhận một lá
mail. Đây cũng là chương mà ý tưởng điều hòa (reconciliation) của
Chương~14 mọc thêm phần mở rộng quan trọng nhất của nó: \emph{operator}.

\section{Operator và custom resource}

Kubernetes có sẵn controller cho các kind dựng sẵn --- Deployment,
Service. Một \emph{operator} là controller cho những kind bạn tự thêm:
một \emph{CustomResourceDefinition} (CRD) dạy API server một danh từ
mới (\code{Kafka}, \code{KafkaTopic}), và vòng lặp điều hòa của
operator biến các thể hiện của nó thành những object bình thường ---
pod, Service, ConfigMap --- rồi giữ chúng hội tụ. Kiến thức vận hành
trước đây nằm trong runbook (cách roll một broker an toàn, cách tính
advertised listener, cách thêm partition) nay được mã hóa trong code
của operator. Strimzi là operator cho Kafka; mẫu hình này là phổ quát
(operator cho Postgres, Prometheus operator, cert-manager).

\section{Cài operator: một việc của admin}

\begin{lstlisting}[language=bash]
# deploy/k8s/install-strimzi.sh (abridged)
kubectl apply --server-side --force-conflicts \
  -f "https://strimzi.io/install/latest?namespace=${NS}" -n "$NS"   (*@\co{1}@*)
kubectl -n "$NS" rollout status deployment/strimzi-cluster-operator
\end{lstlisting}

\begin{description}
  \item[\co{1}] Gói này chứa các CRD, Deployment của operator và RBAC
  của nó --- phần lớn ở \emph{phạm vi cụm} (cluster-scoped). Hai hệ
  quả: nó do admin apply chứ không phải Jenkins (Role của Jenkins,
  Chương~22, dừng ở namespace), và nó được apply \emph{server-side}
  vì các tài liệu CRD quá lớn cho annotation last-applied phía client.
  Lũy đẳng, nên cũng chính là lệnh nâng cấp. Tham số \code{?namespace=}
  giới hạn operator chỉ theo dõi \code{ts}.
\end{description}

Kustomization (Chương~20) liệt kê \code{kafka.yaml}, nhưng file đó
không thể apply cho đến khi script này đã chạy một lần --- thành viên
thứ ba của bộ ``điều kiện tiên quyết do admin làm'', sau namespace và
các secret.

\begin{pitfall}[no matches for kind ``Kafka'' in version v1beta2]
Triệu chứng: apply một manifest Kafka chép từ bất kỳ tutorial nào cũng
thất bại với lỗi này dù operator đã cài và khỏe mạnh. Nguyên nhân:
Strimzi~1.x chỉ phục vụ \code{kafka.strimzi.io/v1}; phiên bản
\code{v1beta2} mà cả thập kỷ bài blog dùng đã biến mất. Hãy kiểm tra
CRD thật sự phục vụ phiên bản nào trước khi tin một ví dụ:
\code{kubectl get crd kafkas.kafka.strimzi.io -o jsonpath=\{.spec.versions[*].name\}}.
Dự án này dính lỗi đó ngay ngày cài đặt.
\end{pitfall}

\section{Cụm, được khai báo}

\begin{lstlisting}[language=yaml]
apiVersion: kafka.strimzi.io/v1
kind: KafkaNodePool                        (*@\co{1}@*)
metadata:
  name: dual-role
  labels: { strimzi.io/cluster: ts-kafka } (*@\co{2}@*)
spec:
  replicas: 1
  roles: [controller, broker]
  storage: { type: persistent-claim, size: 5Gi, deleteClaim: false } (*@\co{3}@*)
  resources:
    requests: { memory: 768Mi, cpu: 200m }
    limits: { memory: 1Gi }
---
apiVersion: kafka.strimzi.io/v1
kind: Kafka
metadata:
  name: ts-kafka
spec:
  kafka:
    listeners:
      - { name: plain, port: 9092, type: internal, tls: false }  (*@\co{4}@*)
    config:
      offsets.topic.replication.factor: 1  (*@\co{5}@*)
      default.replication.factor: 1
      min.insync.replicas: 1
      log.retention.hours: 168
    jvmOptions: { -Xms: 256m, -Xmx: 512m } (*@\co{6}@*)
  entityOperator:
    topicOperator: {}                      (*@\co{7}@*)
\end{lstlisting}

\begin{description}
  \item[\co{1}] Hai object: \emph{pool} nói Kafka chạy ở đâu (replica,
  vai trò, lưu trữ, tài nguyên); resource \code{Kafka} nói nó hành xử
  thế nào. Một node KRaft hai vai trò là cách Kubernetes viết
  \code{process.roles=broker,controller} của Compose.
  \item[\co{2}] Label này nối pool với cụm; operator khớp theo nó. Sai
  label thì im lặng --- pool đơn giản là không bao giờ thành hình.
  \item[\co{3}] Một PVC qua local-path, y hệt các cơ sở dữ liệu
  (Chương~19). \code{deleteClaim: false}: xóa resource Kafka vẫn giữ
  lại log --- cùng sự bất đối xứng mà StatefulSet có theo mặc định.
  \item[\co{4}] So với khối ba listener, hai advertised address trong
  \code{docker-compose.yml}: biến mất. Bạn khai báo \emph{một listener
  nội bộ}; operator tính ra Service bootstrap và Service cho từng
  broker, các advertised address, và listener của KRaft controller.
  Mê cung của Chương~10 là kiến thức vận hành; giờ nó nằm trong
  operator.
  \item[\co{5}] Một broker duy nhất nghĩa là mọi replication factor
  phải bằng 1, nếu không lần tạo topic đầu tiên sẽ thất bại vì đòi
  những broker không tồn tại.
  \item[\co{6}] Heap mặc định của Kafka giả định một máy chủ; nửa GiB
  đủ cho lưu lượng nhỏ giọt của homelab và để dành bộ nhớ cho các
  service JVM.
  \item[\co{7}] Nửa topic của entity operator theo dõi các object
  \code{KafkaTopic} và điều hòa chúng thành topic thật. Nửa user bị
  bỏ qua: không xác thực, không có user nào để quản lý.
\end{description}

Strimzi đặt tên Service bootstrap là \code{<cluster>-kafka-bootstrap};
vì thế mọi producer và consumer nhận
\code{SPRING\_KAFKA\_BOOTSTRAP\_SERVERS=ts-kafka-kafka-bootstrap:9092}
trong Deployment của nó --- mặc định \code{localhost:9092} trong cấu
hình của Phần~III bị ghi đè theo đúng cách mọi địa chỉ khác đã bị ghi
đè.

\begin{decision}[operator, không phải StatefulSet viết tay]
Một StatefulSet KRaft thuần (một pod, một PVC, khối env của compose
chép vào) sẽ tốn ít bộ nhớ hơn --- không operator, không entity
operator --- và đó là phương án dự phòng thẳng thắn nếu tranh chấp
CPU có lúc gây khó. Operator thắng vì đó là thứ các công ty dùng
Kubernetes chạy (Red Hat bán nó dưới tên AMQ Streams), vì topic trở
thành object khai báo, và vì nâng cấp broker cùng khởi động lại cuốn
chiếu trở thành code đã được người khác kiểm thử. Học nó chính là
mục đích; phương án dự phòng vẫn được ghi lại trong
\code{STACK-K3S-JENKINS.md}.
\end{decision}

\section{Topic dưới dạng object}

\begin{lstlisting}[language=yaml]
apiVersion: kafka.strimzi.io/v1
kind: KafkaTopic
metadata:
  name: ts.user.registered
  labels: { strimzi.io/cluster: ts-kafka }
spec: { partitions: 1, replicas: 1 }
\end{lstlisting}

Kafka sẽ tự tạo các topic này ở lần dùng đầu tiên; khai báo chúng
biến số partition thành một quyết định được review, loại bỏ cuộc đua
ngày đầu giữa producer và topic, và biến \code{kubectl get kafkatopics}
thành bản kiểm kê mọi hợp đồng sự kiện trong hệ thống (bốn, khớp với
bảng của Chương~6). Mỗi topic một partition: thứ tự toàn cục, nên
đương nhiên giữ được thứ tự theo từng người dùng mà các publisher
đặt key cho. Partition có thể thêm về sau, không bao giờ bớt được.

\section{Service cuối cùng, và hộp thư của nó}

Manifest của notification-service là Deployment tiêu chuẩn của
Chương~15 với ba điểm riêng: địa chỉ bootstrap ở trên;
\code{SPRING\_MAIL\_HOST=maildev} / \code{PORT=1025} ghi đè mặc định
\code{\$\{MAIL\_HOST:localhost\}} của cấu hình qua relaxed binding;
và các probe trên \code{/actuator/health} --- service này hoàn toàn
không có Spring Security, nên endpoint mở theo mặc định (câu trả lời
thứ ba cho câu hỏi về probe của Chương~17). Maildev là một Deployment
nhỏ xíu kèm Service (cổng 1025 SMTP, 1080 web); nó không gửi đi đâu
cả và hiển thị mọi thứ nó bắt được. Không có Ingress: đây là công cụ
của người vận hành, truy cập qua \code{port-forward}. Nó cũng cung
cấp bài học thực chiến thứ hai của cột mốc này: được apply bằng tay
trước khi pipeline chạy, selector của nó thiếu label mà Kustomize
tiêm vào, và lần deploy đầu tiên của pipeline chết với ``field is
immutable'' --- cạm bẫy của Chương~20, học theo cách đắt giá.

\begin{pitfall}[NotReady --- FatalProblem, trong khi pod đang Running]
Triệu chứng: vài phút sau khi apply, \code{kubectl get kafka} cho thấy
\code{READY} trống và một condition ghi ``Pod is unschedulable or is
not starting'' --- thế nhưng \code{get pods} cho thấy broker
\code{1/1 Running}. Nguyên nhân: vòng điều hòa của operator hết
timeout trong lúc image Kafka 600\,MB vẫn đang tải qua Wi-Fi, và ghi
lại một trạng thái thất bại; vòng lặp định kỳ \emph{tiếp theo} (mỗi
hai phút) quan sát thấy pod đang chạy và chuyển nó sang Ready. Trạng
thái của operator chỉ trung thực về cuối cùng. Trước khi phản ứng với
một condition đỏ, hãy kiểm tra pod mà nó đang mô tả --- và kiểm tra
log của operator, nơi các dòng ``reconciled'' cho thấy vòng lặp vẫn
đang quay.
\end{pitfall}

\begin{hardware}
Operator $+$ broker $+$ entity operator tốn khoảng 1,5\,GB ở đây, với
heap của broker giới hạn 512\,MiB. Lần kéo image đầu tiên là phần đắt
nhất --- hãy chuẩn bị tinh thần mười phút trên Wi-Fi gia đình, và cạm
bẫy trạng thái cũ ở trên. Mọi thứ sau đó đều nhanh: image đã được
cache trên node, điều mà bài kiểm tra khói bên dưới tận dụng.
\end{hardware}

\begin{handson}[chứng minh broker, rồi chứng minh luồng]
Một pod dùng xong vứt với Kafka CLI, ở cùng vị trí mạng với các
service, dùng image node đã có sẵn:
\begin{lstlisting}[language=bash]
kubectl -n ts run kafka-smoke --rm -i --restart=Never \
  --image=quay.io/strimzi/kafka:1.2.0-kafka-4.3.1 -- sh -c '
  B=ts-kafka-kafka-bootstrap:9092
  bin/kafka-topics.sh --bootstrap-server $B --create --topic smoke \
      --partitions 1 --replication-factor 1
  echo hello | bin/kafka-console-producer.sh --bootstrap-server $B --topic smoke
  bin/kafka-console-consumer.sh --bootstrap-server $B --topic smoke \
      --from-beginning --max-messages 1 --timeout-ms 20000
  bin/kafka-topics.sh --bootstrap-server $B --delete --topic smoke'
\end{lstlisting}
Dùng một topic nháp, không phải topic thật: một chuỗi ``hello'' thuần
văn bản trên \code{ts.user.registered} là poison pill cho consumer
JSON (Chương~11) --- được bỏ qua êm thấm nhờ
\code{ErrorHandlingDeserializer}, nhưng việc gì phải phát hiện ra
trong production. Rồi đến luồng thật, khi pipeline đã triển khai xong
các service:
\begin{lstlisting}[language=bash]
curl -X POST http://ts.local/api/auth/register -H 'Content-Type: application/json' \
  -d '{"email":"t@example.com","password":"Passw0rd!","firstName":"Test",
       "lastName":"User","role":"TEACHER","activationMethod":"EMAIL"}'
# -> 201 {"message":"Registration successful. Check your email ..."}
kubectl -n ts logs deploy/notification-service | grep -i registered
# maildev's REST API, from inside the cluster (or port-forward 1080 for the UI):
kubectl -n ts run mailcheck --rm -i --restart=Never \
  --image=curlimages/curl:8.10.1 -- -s http://maildev:1080/email
# -> [{"subject":"Activate your TeacherSupporter account","to":"t@example.com",...
\end{lstlisting}
Request đó đi qua mọi phần của cuốn sách này: Ingress $\rightarrow$
gateway $\rightarrow$ auth-service $\rightarrow$ Postgres $\rightarrow$
Kafka $\rightarrow$ notification-service $\rightarrow$ SMTP. Nó chạy
lần đầu ngày 2026-08-25, pipeline build~\#7 --- sau khi các build
\#3--\#6 mỗi lần dạy một cạm bẫy của chương này và của các Chương~17,
20 và 23. Một đầu mối còn lỏng mà nó phơi ra: link kích hoạt trong
mail vẫn ghi \code{localhost:8080}, mặc định dev của URL gốc cho link
--- một ghi đè env dành cho cột mốc~7, khi frontend cho cụm một địa
chỉ công khai thật sự.
\end{handson}

\section{Ngoài dự án}

Kafka được quản lý --- AWS MSK, Confluent Cloud, endpoint Kafka của
Azure Event Hubs --- là operator của operator: cùng một khai báo hình
dạng \code{Kafka} nhưng thực hiện trong console của cloud, với những
broker bạn không bao giờ thấy. Chuyển từ đây sang đó là một chuỗi
bootstrap cộng SASL/TLS trên listener; các object \code{KafkaTopic}
sẽ nhường chỗ cho resource Terraform nói đúng điều tương tự. Bài học
còn lại để mang theo: mẫu hình operator là cách Kubernetes lưu trú
\emph{mọi} thành phần nền tảng có trạng thái --- Postgres
(CloudNativePG), Redis, Elasticsearch --- và khi đã cài, đã dùng sai
phiên bản, đã debug trạng thái của một cái, bạn biết hình dạng của
tất cả chúng.

\begin{quiz}
\begin{enumerate}
  \item Định nghĩa operator theo cách nói về controller của Chương~14,
  và nêu hai kind object mà Strimzi đã thêm vào cụm này.
  \item Vì sao \code{install-strimzi.sh} là bước của admin trong khi
  \code{kafka.yaml} do Jenkins apply? Khái niệm RBAC nào vạch ra ranh
  giới đó?
  \item Liệt kê ba thứ mà khối Kafka của Compose cấu hình bằng tay
  nhưng resource \code{Kafka} không hề nhắc tới, và nói ai tính chúng
  bây giờ.
  \item Bạn thấy \code{NotReady / FatalProblem} trên resource Kafka.
  Hai thứ bạn kiểm tra trước khi hành động là gì, và vì sao trạng thái
  đó có thể sai?
  \item Vì sao khai báo topic dưới dạng object khi Kafka tự tạo được
  chúng, và vì sao mỗi topic chỉ một partition?
\end{enumerate}
\end{quiz}
